\section{Introduction}\label{sec:intro}

Programming courses assign many small exercises, and learners progress at very different speeds. An exercise that is too difficult discourages a learner; one that is too easy teaches little. Adaptive exercise recommendation addresses this by estimating each learner's state from past attempts and choosing the next exercise accordingly, an idea often motivated by the zone of proximal development \citep{vygotsky1978mind} and with a long tradition in adaptive educational systems for programming \citep{brusilovsky2007user,hsiao2010guiding}. A now common recipe combines a knowledge-tracing (KT) model, which forecasts a learner's success on each exercise from their history, with a reinforcement-learning (RL) agent that treats the choice of exercise as a sequential decision problem: the learner model's state is the agent's state, recommending an exercise is an action, and an increase in the model's predicted knowledge is the reward \citep{huang2019dre,liu2019cseal,ai2019cdkt,chen2023src,li2023gehrl,ozyurt2025exrec}. Because trials with learners are expensive, such agents are often trained \emph{and evaluated} inside a fitted learner model, and the resulting gain in predicted knowledge is reported as the benefit of the policy \citep{cheng2026unier}.

This evaluation is only as good as the simulator's account of how practice changes success on other exercises. KT models are fitted and validated on next-step forecasts of the sequences that learners actually followed. A policy, however, creates new sequences, and the simulator must extrapolate what these sequences would do. Two models can forecast observed data equally well and still disagree about such counterfactuals \citep{breiman2001two,damour2022underspecification}. For instructional decisions this was noticed early: individualised and population knowledge-tracing parameters imply very different amounts of practice \citep{lee2012individualizing}, and student models with similar predictive accuracy can recommend substantially different amounts of practice \citep{rollinson2015predictive}. The robust evaluation matrix therefore evaluates candidate policies under several student models \citep{doroudi2017rem}, and a review of RL for instructional sequencing found that policies derived from purely data-driven models often failed to outperform reasonable baselines when tested with students \citep{doroudi2019reward}. In reinforcement learning more generally, an optimiser trained against a learned model or a learned reward exploits that model's errors \citep{kurutach2018metrpo,janner2019trust,gao2023scaling}. Modern KT--RL pipelines combine both risks---the learned model is at once the environment, the reward and the yardstick---yet it is rarely quantified how much of a reported gain reflects learning that the simulator encodes and how much reflects the simulator's idiosyncrasies.

We study this question for introductory programming, using {{n_submissions}} judged C/C++ submissions by {{n_learners}} learners of the ACcoding online judge \citep{chen2024accoding}. We fit six learner models of equal predictive accuracy, including deep knowledge tracing (DKT) as the deployed tracker and an entropy-based identification tree for interpretable diagnosis; train Double DQN and PPO recommenders in three kinds of training worlds; and evaluate every policy in nine worlds that differ in the learners they were fitted to and in how learning is represented. Two instruments make the evaluation diagnostic. A \emph{negative-control} world without any learning, in which the true gain of every policy is exactly zero, exposes spurious gains directly. An exact \emph{decomposition} of a tracker's reported gain into the model's own expected forecast change (\emph{drift}) and evidence \emph{surprise} shows where such gains come from. We prove that drift vanishes for coherent Bayesian models without learning and that, inside its own simulator, the expected gain of any policy equals the drift it accumulates; in-simulator evaluation therefore rewards a policy for choosing problems that the model \emph{expects} to raise its own forecast.

We address four research questions. RQ1: How well do cognitive-state models forecast first-attempt success of unseen programming learners, and do they differ? RQ2: How large are the gains of RL recommenders in the KT simulator they are trained in? RQ3: Do these gains survive independently fitted worlds, and how large are the gains a tracker reports when no learning is possible? RQ4: What produces the discrepancy, and which policies remain robust? All protocols, including four hypotheses, were locked on a development cohort before the evaluation population was loaded.

Our contributions are: (1) an end-to-end, reproducible pipeline for adaptive programming exercise recommendation on real data, from learner models to trained DQN and PPO agents and an interactive application, with pre-registered hypotheses tested on a population never used before; (2) evidence that, although six learner models forecast equally well, the in-simulator gains of KT-reward agents shrink to near zero in independently fitted worlds with latent learning and are almost perfectly rank-correlated with the gains the tracker reports in a world where no learning occurs; (3) the drift decomposition and two propositions that explain these gains as the tracker's accumulated self-predicted change rather than learning, a model-internal diagnostic that needs no additional data; and (4) exploratory evidence that agents trained in a world with explicit learning structure transfer best across latent worlds of that family---but not to neural worlds---together with inexpensive checks for simulator-based evaluation.

\Cref{sec:related} reviews related work, \cref{sec:data} describes the data and \cref{sec:methods} the models, decision problem, evaluation worlds and analysis. \Cref{sec:results} reports the results, \cref{sec:discussion} discusses their implications and limitations, and \cref{sec:conclusion} concludes.
